\section{Continuous Assessment 2 (CA-2) --- Set B}
\label{sec:ca2_set_b}

\LPUPaperHeader{Continuous Assessment (CA-2)}{B}{60 Minutes}{30 Marks}{CO4 (Environmental Pollution \& Waste), CO5 (Disaster Management)}

\begin{center}
\sffamily\textbf{\color{AlertRed}Note: Attempt all six questions. Each question carries 5 marks.}
\end{center}

\questionitem{1}{
    Evaluate the National Ambient Air Quality Standards (NAAQS) for $\text{PM}_{10}$ and $\text{PM}_{2.5}$. Explain why fine particulate matter ($\text{PM}_{2.5}$) poses severe cardiovascular and pulmonary health risks compared to coarse particulates.
}{5}{CO4}{L3 Apply}

\questionitem{2}{
    Compare the separation mechanisms and operational limits of Cyclone Separators and Fabric Filter Baghouses for industrial flue gas cleaning.
}{5}{CO4}{L3 Apply}

\questionitem{3}{
    Explain the biochemical mechanism of aquatic Eutrophication triggered by agricultural nitrogen and phosphorus runoff, detailing the progression from algal bloom to benthic anoxia.
}{5}{CO4}{L2 Understand}

\questionitem{4}{
    Diagnose the technological and human failures responsible for the 1984 Bhopal Gas Tragedy. List three safety systems that were non-operational during the runaway reaction.
}{5}{CO5}{L4 Analyze}

\questionitem{5}{
    Detail the institutional hierarchy established under the Disaster Management Act, 2005 (NDMA, SDMA, DDMA, NDRF) and identify their respective executive heads.
}{5}{CO5}{L2 Understand}

\questionitem{6}{
    Examine the four Priorities for Action of the Sendai Framework for Disaster Risk Reduction (2015--2030) and explain the concept of ``Build Back Better''.
}{5}{CO5}{L3 Apply}

\newpage
\subsection*{Solutions and Marking Rubrics --- CA-2 Set B}

\begin{solutionbox}[Solution to Question 4 --- Bhopal Diagnostic]
\textbf{1. Overview \& Chemical Agent (1.5 Marks):}\\
On Dec 2--3, 1984, over $42\text{ tonnes}$ of Methyl Isocyanate (MIC) gas leaked from Tank E610 at the Union Carbide plant in Bhopal due to exothermic reaction following water ingress.\\
\textbf{2. Decommissioned Safety Systems (2.5 Marks):}\\
(a) The refrigeration cooling system was turned off to reduce electricity expenses.\\
(b) The caustic soda vent gas scrubber was non-functional.\\
(c) The flare tower piping was dismantled for maintenance without backup.\\
\textbf{3. Environmental \& Statutory Impact (1 Mark):}\\
Induced thermal inversion held the gas at ground level, killing thousands and leading to the enactment of the Environment (Protection) Act, 1986.
\end{solutionbox}
\newpage
